%%%%%%%%%%%%%%%%%%%% VIBETEST QWEN KAGGLE -- CLEAN SUBMISSION DRAFT %%%%%%%%%
%
% This is the submission-facing draft. The ACL style file controls margins and
% layout; do not add geometry/fullpage packages, since they will break the
% conference margin requirements.
%
%%%%%%%%%%%%%%%%%%%%%%%%%%%%%%%%%%%%%%%%%%%%%%%%%%%%%%%%%%%%%%%%%%%%%%%%%%%%%%%

\documentclass[11pt]{article}

\usepackage[review]{acl}

\usepackage{times}
\usepackage{latexsym}
\usepackage[T1]{fontenc}
\usepackage[utf8]{inputenc}
\usepackage{microtype}
\usepackage{booktabs}
\usepackage{graphicx}
\graphicspath{{figures/}}
\usepackage{multirow}
\usepackage{xcolor}
\usepackage{amsmath}
\usepackage{array}
\usepackage{tabularx}
\usepackage{url}
\usepackage{hyperref}
\usepackage[capitalize,noabbrev]{cleveref}

\definecolor{darkblue}{rgb}{0,0,0.5}
\hypersetup{colorlinks=true, citecolor=darkblue, linkcolor=darkblue, urlcolor=darkblue}

% REVIEW-ONLY: flags spots still needing author input/decision. Renders in red.
% Robust so it is safe inside captions and other moving arguments.
% Before camera-ready, hide all of them with:
%   \DeclareRobustCommand{\needsinput}[1]{}
\DeclareRobustCommand{\needsinput}[1]{\textcolor{red}{\textbf{[INPUT NEEDED:}~#1\textbf{]}}}

\newcommand{\pass}{\textsc{pass}}
\newcommand{\fail}{\textsc{fail}}
\newcommand{\inconclusive}{\textsc{inconclusive}}
\newcommand{\vibetest}{\textsc{VibeTest}}
\newcommand{\qwen}{Qwen3.6}

\title{\vibetest: Evidence-Guided Agentic Testing\\for Natural-Language ML Pipeline Properties}

\author{Anonymous ACL submission}

\begin{document}
\maketitle

\begin{abstract}
ML repositories are governed by requirements that are easy to state but hard to
test: avoid leakage, evaluate on held-out data, apply augmentation only to
training examples, or compute metrics according to their definitions.
These requirements are natural-language properties over messy repository
context and are rarely written down as formal assertions.
We propose \vibetest, an agentic testing framework that validates a repository
against a natural-language ML-pipeline specification, without requiring a formal
one, and returns a \pass, \fail, or \inconclusive{} verdict with an evidence
report and a continuous confidence score.
On 1,125 synthetic Kaggle (repository, property) pairs, \vibetest{} with \qwen{}
achieves 0.802--0.837 selective macro F1 across prompt-example settings and
outperforms reviewer-style baselines.
On 1,920 real Kaggle cases, a conservative sampled fail audit estimates
0.931--0.938 macro F1.
The method also transfers across backbones: GPT-5-mini reaches comparable macro
F1 at similar token cost.
But the residual errors carry a second finding.
When we hand-audit 53 high-confidence cases where the model and the benchmark
disagree, only 29 are genuine model misses; the other 24 are cases where the
injected label and the evidence disagree about what the property means.
Selective F1 alone therefore conflates model quality with label
operationalization, and we argue that natural-language repository testing should
report both.
\end{abstract}

\section{Introduction}

Modern ML systems are usually delivered as full repositories.
Their behavior depends on choices scattered across notebooks, scripts, data
loaders, training loops, and evaluation code.
Many properties practitioners care about are straightforward to describe:
use the held-out split, avoid train-test leakage, compute ROC-AUC from scores
rather than hard labels, or ensure that intended parameters are updated during
training.
Yet these properties are hard to test automatically because they are expressed
in natural language and must be grounded in repository evidence.

This setting sits between code review and formal verification.
Program analysis can provide strong guarantees once a property has been
formalized~\citep{cousot1979systematic,nielson1999principles}, and query
systems such as CodeQL are effective when developers can encode a semantic
pattern~\citep{codeql}.
ML debugging systems catch important training failures through structured
checks~\citep{schoop2021umlaut,traincheck,kang2020assertions,naik2024torchql}.
But many ML-pipeline requirements first appear as informal intent.
For these cases, running the check is the easy part.
The problem is deciding what evidence in a repository counts as satisfying or
violating the property.

LLM agents are a natural fit for this evidence-gathering step.
They can inspect files, follow code paths, read notebooks, and explain their
decisions.
The same flexibility makes evaluation difficult.
A model may report a failure because it found real evidence, because it noticed
a best-practice concern, or because the repository lacks enough information to
decide.
These outcomes should not be collapsed.
For example, the absence of a diagnostic sanity check is not proof that the
repository violates a property; it may only make the property unverifiable.

We introduce \vibetest, an agentic testing framework for natural-language
ML-pipeline properties, as a way to close this gap.
Given a repository $R$ and property $p$, \vibetest{} emits one of three
verdicts: \pass, \fail, or \inconclusive{} (\Cref{fig:overview}).
It also emits an evidence report and a continuous \texttt{case\_score}, where low
values indicate confidence in \pass{} and high values indicate confidence in
\fail.
This score enables selective evaluation: high-confidence cases are accepted and
ambiguous cases are abstained.

Two findings emerge.
Evidence-guided testing is already useful: on synthetic Kaggle-style
repositories \vibetest{} outperforms reviewer-style baselines, and on real
Kaggle notebooks a conservative sampled audit puts macro F1 near 0.93.
More interesting is what the residual errors are made of.
When we audit the high-confidence cases where \qwen{} and the synthetic label
disagree, nearly half turn out to be disagreements about where a property's
boundary lies: the injected change, the benchmark label, the prompt, and the
auditor can each place it somewhere different.
Testing a natural-language property, then, includes deciding what the property
means.

\begin{figure*}[t]
  \centering
  \includegraphics[width=0.95\linewidth]{agentic_testing_evidence_verification.png}
  \caption{Evidence-guided agentic testing. A repository and a natural-language
  ML-pipeline specification are given to the testing agent. The agent gathers
  repository evidence and emits a verdict, a confidence score, and an evidence
  report. Downstream audits check whether predicted failures are supported by
  the reported repository context.}
  \label{fig:overview}
\end{figure*}

We make four contributions.
First, we formulate natural-language ML-pipeline checking as evidence-guided
repository testing with explicit \pass/\fail/\inconclusive{} verdicts.
Second, we propose \vibetest, an agentic testing framework that produces these
verdicts with cited evidence and a calibratable confidence score.
Third, we evaluate \vibetest{} on synthetic and real Kaggle benchmarks with two
backbones (\qwen{} and GPT-5-mini), comparing against reviewer-style baselines
and direct-prompting probes, and measuring whether the cited evidence supports
each verdict.
Fourth, we audit synthetic disagreements and show that raw synthetic F1 mixes
model misses with property-definition mismatches.

\section{Problem Setting}
\label{sec:problem}

We evaluate a (repository, property) pair $(R,p)$, where $R$ is an ML repository
and $p$ is a natural-language specification of expected behavior, such as ``no
training on the test set.''
The task is to return a decision $y \in \{\pass, \fail, \inconclusive\}$ together
with an evidence report $E$ that identifies the files, executions, and outputs
justifying $y$ and lets a reader verify it independently.
This differs from generic code review, where any plausible defect may be
reported, and from formal verification, where the property is already expressed
in a formal language.

\paragraph{Informal properties.}
Natural-language properties have boundaries.
``Avoid replaceable Python loops'' might target elementwise numerical loops in
tensor code while excluding loops that only print file paths.
``Do not freeze trainable parameters unless intended'' requires deciding what
counts as explicit intent.
The tester must operationalize the property before it can evaluate code.

\paragraph{Incomplete evidence.}
ML repositories often omit artifacts needed for a definitive judgment.
A notebook may not run, external data may be unavailable, or a metric may be
reported without the exact training artifact.
In these cases, a forced binary decision can turn missing evidence into false
failures.
\inconclusive{} must therefore be a first-class outcome.

\paragraph{Benchmark labels as operationalizations.}
Synthetic labels come from injected changes, and an injected change is only one
operationalization of the natural-language property.
When a model and label disagree, either may be wrong, or they may be applying
different but defensible property boundaries.
Evidence audits are therefore part of the benchmark itself.

\section{\vibetest}
\label{sec:method}

We design \vibetest{} as a repository-level testing agent for natural-language
ML properties.
Its guiding principle is evidence discipline: a \fail{} verdict must be supported
by specific repository evidence, and weak or missing evidence yields
\inconclusive{}.
\vibetest{} operates on a sandboxed copy of the repository.
It may run and modify code to probe a property, for example by adding logging or
perturbing labels, but these changes do not persist to the original repository.

\paragraph{Output schema.}
For each case, \vibetest{} returns a verdict, a \texttt{case\_score}, and an
evidence report: concise reasoning together with the cited files, outputs, and
logs that justify the verdict.
\pass{} means the repository appears to satisfy the property with supporting
evidence.
\fail{} means the repository appears to violate the property with supporting
evidence.
\inconclusive{} means the repository does not provide enough evidence to
decide.
\Cref{fig:worked-example} shows a representative report.

\begin{figure*}[t]
  \centering
  \small
  \fbox{\begin{minipage}{0.80\linewidth}
  \textbf{Property.} Any data augmentation is applied only to the training data;
  validation and test use deterministic preprocessing.\\[2pt]
  \textbf{Verdict.}~\fail\qquad \texttt{case\_score}~$=0.95$\\[2pt]
  \textbf{Reasoning.} Noise injection is applied to the full feature tensor in
  Cell~45, \emph{before} the train/validation split in Cell~74, so the
  validation split inherits the augmentation. Only the test set (Cell~114)
  stays deterministic.\\[3pt]
  \textbf{Cited evidence.}
  \begin{itemize}\setlength{\itemsep}{1pt}\setlength{\topsep}{1pt}
    \item Cell~45 adds noise to the whole tensor:\\
      {\ttfamily\footnotesize t\_indep[:,jitter\_cols] = t\_indep[:,jitter\_cols]
      + torch.randn\_like(\ldots)*0.02}
    \item Cell~74 then splits the already-augmented tensor:\\
      {\ttfamily\footnotesize trn\_indep, val\_indep = t\_indep[trn\_split],
      t\_indep[val\_split]}
    \item Cell~114 preprocesses the test set with deterministic operations only.
  \end{itemize}
  \end{minipage}}
  \caption{A \vibetest{} report on a synthetic Titanic repository (\qwen).
  The injected bug augments features before the train/validation
  split, contaminating validation while leaving the test set clean; \vibetest{}
  returns \fail{} with a high \texttt{case\_score} and cites the three cells that
  establish the violation. Reasoning is lightly condensed for space.}
  \label{fig:worked-example}
\end{figure*}

\paragraph{Selective scoring.}
The \texttt{case\_score} lies in $[0,1]$.
Scores near 0 indicate confidence that the property is satisfied, while scores
near 1 indicate confidence that the property is violated.
For two-sided thresholding, a threshold $t$ accepts \pass{} when
\texttt{case\_score} $\leq 1-t$, accepts \fail{} when
\texttt{case\_score} $\geq t$, and abstains otherwise.
Coverage is the fraction of cases accepted as \pass{} or \fail.

\paragraph{Prompt-example settings.}
We evaluate three static prompt settings.
\textbf{ex0} uses no representative examples.
\textbf{ex10} includes ten concise examples of representative failure cases.
\textbf{ex20} includes twenty such examples.
The examples are intended to calibrate the evidence standard; they do not
provide answers for the evaluated (repository, property) pairs.

\paragraph{Connection to structured reasoning.}
Prior work on verifiers, process supervision, and self-refinement shows that
intermediate structure can make model outputs more reliable
~\citep{cobbe2021verifier,lightman2024process,madaan2023selfrefine,weng2023selfverification}.
PIPS applies this idea to per-instance program synthesis by using structural
feedback~\citep{stein2025pips}.
\vibetest{} applies the same high-level lesson to repository testing: the
intermediate object is an evidence-backed verdict that can be audited.

\section{Experimental Setup}
\label{sec:setup}

\paragraph{Synthetic Kaggle.}
The synthetic benchmark contains 75 Kaggle-style notebooks across three
competitions of different modalities: Titanic (tabular), Disaster Tweets (text),
and Diabetic Retinopathy (image).
Each notebook is evaluated against 15 ML-pipeline properties, yielding 1,125
(repository, property) pairs.
Labels come from controlled bug injection and are treated as the canonical
synthetic ground truth before audit.

\paragraph{Real Kaggle.}
The real benchmark contains 121 public Kaggle notebooks: 50 Titanic, 50
Disaster Tweets, and 21 Diabetic Retinopathy notebooks
~\citep{kaggle_titanic,kaggle_disaster_tweets,kaggle_diabetic_retinopathy}.
Each notebook is evaluated against up to 16 properties, yielding 1,920 cases.
Because exhaustive real ground truth is unavailable, headline real results use
a conservative sampled audit of predicted failures.
A smaller full-context reaudit is reported only as sensitivity analysis.

\paragraph{Baselines and probes.}
We compare against reviewer-style baselines and a Codex reviewer on the full
synthetic benchmark, evaluated against the same synthetic labels.
Codex is included as an execution/format-fragile baseline: at scale it emits a
parseable committed verdict on only 5.5\% of cases (\Cref{tab:headline}).
We additionally report matched direct-property and Codex-style probes on a
five-repository Titanic subset (\Cref{tab:probe}).
These probes help separate the value of the evidence rubric from simply asking
a model about the property.
To isolate the value of evidence-guided testing from the underlying model, we
also run a \emph{same-checkpoint} probe: on the five-repository Titanic subset we
compare direct-property prompting and \vibetest{} on an identical
Qwen3.6-35B-A3B checkpoint served through OpenRouter (2026-04-15).
Direct prompting sees the notebook and the property and answers directly, with
no evidence rubric, tools, or execution.
On this matched checkpoint, \vibetest{} reaches 0.811 macro F1 and 0.784 failure
recall versus 0.704 and 0.667 for direct prompting, so the gain is attributable
to the evidence-guided procedure.
A full-benchmark same-checkpoint run remains future work.
On real Kaggle we do not report Codex or direct-property numbers, because our
real ground truth comes from fail audits of \vibetest{}'s own predictions and no
comparable audit exists for these baselines.
\needsinput{confirm this framing is acceptable, or run a matched real-Kaggle fail
audit for Codex/direct so they can appear in the real results.}

\paragraph{Metrics.}
For synthetic data, where we have injected ground truth, we report selective
macro F1 over covered cases together with coverage.
Real data has no exhaustive ground truth, so the real macro F1 is only partly
audited: we verify the \fail{} class by sampling predicted failures and checking
whether the cited evidence supports them, but the \pass{} class is taken from
the model's own verdicts.
The estimate therefore leans on failure precision and does not establish recall
over real defects the model missed; we treat it conservatively for that reason,
and read it alongside coverage from the full model outputs.
We also report score histograms and worst-class selective error, defined as the
maximum of the accepted-fail false discovery rate and accepted-pass false
omission rate.
For selective macro F1 we report 95\% confidence intervals from a case-level
bootstrap that resamples the covered (repository, property) pairs with
replacement (20{,}000 replicates); for audited proportions we report Wilson
score intervals.

\paragraph{Audit protocol.}
Both audits are adjudicated by an LLM verifier.
The verifier is a different model family from the tester: the tester is
\qwen{}, while the verifier is GPT-5.4-mini for the synthetic evidence checks
and GPT-5-mini for the real fail audit, both distinct from the tester.
It is shown the property, the tester's verdict, and the cited evidence, and it
decides whether the evidence supports the verdict.
The same verifier provides the per-verdict evidence-support signal we report as
evidence quality (RQ4).
Using a separate model to adjudicate gives a measure of cross-model
independence: a disagreement flagged as a property-operationalization mismatch
is one where the injected label conflicts with an evidence-based reading from a
third model, separate from both the label injector and the tester.
The verifier is itself an LLM and applies the same evidence standard as the
tester; we return to this in \Cref{sec:limitations}.

\section{Results}
\label{sec:results}

\paragraph{RQ1: Does evidence-guided testing improve synthetic performance?}
Yes, and by a clear margin.
At the selective-F1-optimal operating point (\Cref{tab:headline}; the apex of
each risk--coverage curve in \Cref{fig:synthetic-f1}), \vibetest{} reaches
0.802 macro F1 for ex0, 0.813 for ex10, and 0.837 for ex20, against 0.669--0.722
for the reviewer baselines (95\% bootstrap CIs in \Cref{tab:headline}).
The ex10 and ex20 intervals sit clearly above the reviewer intervals, while
ex0's margin over the best reviewer is small.
These are case-level intervals; a repository-level bootstrap would be wider
(\Cref{sec:limitations}), so we read the gain over reviewers as clear for the
example settings but not decisive for ex0.
The ex0, ex10, and ex20 intervals overlap one another, so we do not treat the
example count as a reliable ranking.
This gap is not an artifact of \vibetest{} abstaining more: because the
reviewers emit no abstention score, \Cref{fig:synthetic-f1} charitably folds
their \inconclusive{} verdicts into \pass{} and sweeps a pseudo-score, yet even
under this generous treatment they never reach \vibetest{}, and \vibetest{}
dominates across the shared coverage region.
Left to its native verdicts, a reviewer commits to a \pass{} or \fail{} on only
0.26--0.32 of cases, so the folding credits it with many unjudged passes.
Codex reviewer and TrainCheck are execution/format-fragile at scale: on the full
1{,}125-case synthetic benchmark Codex emits a parseable committed verdict on
only 5.5\% of cases (macro F1 0.419 on that sliver, mostly by predicting \fail),
and TrainCheck produced no usable coverage in our canonical execution.
We report both at their native operating points (\Cref{tab:headline},
\Cref{fig:synthetic-f1}).
These numbers reflect execution fragility in our setting; either tool may
perform better with cleaner integration.

\begin{table*}[t]
  \centering
  \small
  \begin{tabular}{llrrrr c l}
    \toprule
    Benchmark & Method & Examples & Tests & Coverage & Macro F1 & 95\% CI & Label source \\
    \midrule
    Synthetic & \vibetest{} static & 0 & 1125 & 0.560 & 0.802 & {\scriptsize[0.77, 0.83]} & Synthetic GT \\
    Synthetic & \vibetest{} static & 10 & 1125 & 0.643 & 0.813 & {\scriptsize[0.78, 0.84]} & Synthetic GT \\
    Synthetic & \vibetest{} static & 20 & 1125 & 0.562 & 0.837 & {\scriptsize[0.81, 0.87]} & Synthetic GT \\
    Synthetic & Reviewer mode 0 & -- & 1125 & 0.973 & 0.669 & {\scriptsize[0.64, 0.70]} & Synthetic GT \\
    Synthetic & Reviewer mode 1 & -- & 1125 & 0.963 & 0.722 & {\scriptsize[0.69, 0.75]} & Synthetic GT \\
    Synthetic & Reviewer mode 2 & -- & 1125 & 0.982 & 0.711 & {\scriptsize[0.68, 0.74]} & Synthetic GT \\
    Synthetic & Codex reviewer & -- & 1125 & 0.055 & 0.419 & {\scriptsize[0.33, 0.54]} & Synthetic GT \\
    Synthetic & TrainCheck & -- & 1125 & 0.000 & -- & -- & Synthetic GT \\
    Real & \vibetest{} static & 0 & 1920 & 0.822 & 0.931 & -- & Conservative audit \\
    Real & \vibetest{} static & 10 & 1920 & 0.869 & 0.933 & -- & Conservative audit \\
    Real & \vibetest{} static & 20 & 1920 & 0.871 & 0.938 & -- & Conservative audit \\
    \bottomrule
  \end{tabular}
  \caption{Headline results. For synthetic rows, coverage and macro F1 are
  reported at a single operating point (the apex of each method's
  risk--coverage curve in \Cref{fig:synthetic-f1}), so the two columns describe
  the same threshold. \vibetest{} peaks at 0.56--0.64 coverage; its full-output
  coverage without extra thresholding is 0.77--0.90. Reviewer baselines emit no
  abstention score, so \Cref{fig:synthetic-f1} folds their \inconclusive{}
  verdicts into \pass{} and sweeps a pseudo-score (charitable to the baseline,
  which is then credited with unjudged passes); their native
  non-\inconclusive{} rate is only 0.26--0.32. Codex reviewer and TrainCheck are
  execution/format-fragile at scale: on the full benchmark they emit a parseable
  committed verdict on only 5.5\% and 0\% of cases respectively, so they are
  reported at their native committed operating point without folding. (On the
  small five-repository probe of \Cref{tab:probe}, Codex was scored with the
  reviewer folding and so appears at higher coverage there.
  \needsinput{Codex is 0.055 coverage here vs 0.44 in the probe table; decide
  whether to keep both with this reconciling note or drop the probe Codex column
  to avoid the apparent conflict.}) Real rows report the
  conservative sampled fail audit; real coverage is the non-\inconclusive{} rate
  over all 1{,}920 outputs, which uses a different sample than the audit
  (\Cref{sec:setup}). The 95\% CI is a case-level bootstrap over covered
  (repository, property) pairs; real rows are audit estimates without a
  comparable interval (their auditable fail precision is 0.80 [0.73, 0.86],
  \Cref{sec:real-examples}).}
  \label{tab:headline}
\end{table*}

\begin{figure}[t]
  \centering
  \includegraphics[width=\linewidth]{paper_synthetic_selective_f1_coverage.pdf}
  \caption{Selective macro F1 versus coverage on synthetic Kaggle. \vibetest{}
  occupies the higher-F1 region and dominates the reviewer baselines wherever
  their coverage overlaps. Reviewer baselines have no abstention score, so their
  \inconclusive{} verdicts are folded into \pass{} before sweeping a pseudo-score;
  this inflates their coverage and is charitable to the baseline. Codex reviewer
  (full run) is a single point at the bottom left: it commits on only 5.5\% of
  cases at scale. TrainCheck has zero usable coverage and appears only in
  \Cref{tab:headline}.}
  \label{fig:synthetic-f1}
\end{figure}

\paragraph{Prompt-only probes.}
\Cref{tab:probe} compares \vibetest{} to direct-property and Codex-style probes
on the same five Titanic synthetic repositories.
The Qwen-family direct-property probe reaches 0.486 failure recall, while
\vibetest{} ex10 reaches 0.811 on the same subset.
This suggests that representative examples and the evidence rubric improve
failure discovery beyond making the model more verbose.

\begin{table}[t]
  \centering
  \scriptsize
  \setlength{\tabcolsep}{3pt}
  \begin{tabular}{@{}>{\raggedright\arraybackslash}p{0.42\linewidth}rrrr@{}}
    \toprule
    Method & Cov. & F1 & P & R \\
    \midrule
    Direct Qwen Flash & 0.840 & 0.620 & 0.750 & 0.486 \\
    Direct GPT-4.1 mini & 0.973 & 0.610 & 0.824 & 0.378 \\
    Codex reviewer & 0.440 & 0.349 & 0.625 & 0.135 \\
    \vibetest{} ex0 & 0.840 & 0.681 & 0.840 & 0.568 \\
    \vibetest{} ex10 & 0.920 & 0.791 & 0.833 & 0.811 \\
    \vibetest{} ex20 & 0.893 & 0.775 & 0.794 & 0.730 \\
    \bottomrule
  \end{tabular}
  \caption{Matched five-repository Titanic probe with 75 cases. P/R denote
  failure precision and recall. Direct prompting and generic review have weaker
  failure recall than \vibetest{} with prompt examples.}
  \label{tab:probe}
\end{table}

\paragraph{RQ2: Are scores useful for selective prediction?}
The score supports useful selective curves, but it is not perfectly calibrated.
\Cref{fig:synthetic-error} shows that worst-class selective error remains
nontrivial even when abstention removes many cases, and
\Cref{fig:score-hists} shows the underlying score distributions: the model
concentrates mass near the confident \pass{} and \fail{} ends but still places
appreciable mass in the ambiguous middle.
This motivates auditing the disagreements directly, which we do in
\Cref{sec:audit}.

\begin{figure}[t]
  \centering
  \includegraphics[width=\linewidth]{paper_synthetic_max_error_coverage.pdf}
  \caption{Synthetic worst-class selective error versus coverage on a linear
  scale. Lower is better. Residual error remains visible after abstention,
  motivating the disagreement audit in \Cref{sec:audit}. Coverage below 0.15 is
  omitted: with so few accepted cases, a single misclassification swings the
  worst-class rate and the estimate is not meaningful.}
  \label{fig:synthetic-error}
\end{figure}

\begin{figure*}[t]
  \centering
  \includegraphics[width=\linewidth]{paper_case_score_ecdf.pdf}
  \caption{Empirical case-score distributions for synthetic Kaggle (left) and
  real Kaggle (right). Steep segments indicate score clusters; flatter segments
  indicate score ranges the model uses less often. The curves show score spread
  without letting high-count bins dominate the visualization.}
  \label{fig:score-hists}
\end{figure*}

\paragraph{RQ3: Does the pattern transfer to real Kaggle?}
Conservative real-audit estimates are strong across example settings.
\Cref{fig:real-f1} shows estimated macro F1 around 0.93, with ex20 slightly
strongest in the headline table.
Full-context reaudit suggests that some conservative false-fail penalties were
too harsh; we treat that analysis only as a sensitivity check and keep the
conservative values as the headline.

\begin{figure}[t]
  \centering
  \includegraphics[width=\linewidth]{paper_real_conservative_f1_coverage.pdf}
  \caption{Conservative macro F1 versus coverage on real Kaggle. The sampled
  fail audit estimates roughly 0.93 macro F1 across prompt-example settings over
  the reported coverage range. Coverage below 0.30 is omitted because too few
  accepted failures enter the audit estimate there for it to be stable; the
  headline uses the high-coverage region.}
  \label{fig:real-f1}
\end{figure}

\begin{table}[t]
  \centering
  \small
  \begin{tabular}{lr}
    \toprule
    Examples & Conservative macro F1 \\
    \midrule
    0 & 0.931 \\
    10 & 0.933 \\
    20 & 0.938 \\
    \bottomrule
  \end{tabular}
  \caption{Real Kaggle headline (conservative sampled fail audit). As a
  one-sided sensitivity check, re-auditing the conservative false failures with
  full source context accepts 26 of 30 as supported, which would push macro F1
  into the 0.98--0.99 range. We treat that only as an optimistic upper bound,
  not an independent estimate, and keep the conservative values as the headline.}
  \label{tab:real-sensitivity}
\end{table}

\paragraph{Does the result depend on the model or the dataset?}
To check that evidence-guided testing is not specific to \qwen{}, we also run
the tester with GPT-5-mini.
GPT-5-mini emits verdicts but no \texttt{case\_score}, so we compare both models
at their native operating point (covered cases are the non-\inconclusive{}
verdicts), which is why the \qwen{} macro F1 here is below its selective peak in
\Cref{tab:headline}.
\Cref{tab:models} shows the two backbones are close in macro F1 (mean 0.769 for
\qwen{}, 0.746 for GPT-5-mini) at comparable coverage and output-token cost
(about 3k tokens per property each).
Both models are weakest on Diabetic Retinopathy (image) and strongest on
Disaster Tweets (text), so difficulty tracks the dataset more than the backbone.

\begin{table}[t]
  \centering
  \small
  \begin{tabular}{llrrr}
    \toprule
    Dataset & Model & Cov. & Macro F1 & Tok/prop \\
    \midrule
    \multirow{2}{*}{Titanic} & \qwen{} & 0.720 & 0.762 & 2.8k \\
     & GPT-5-mini & 0.795 & 0.765 & 3.5k \\
    \midrule
    \multirow{2}{*}{Disaster Tweets} & \qwen{} & 0.835 & 0.795 & 3.1k \\
     & GPT-5-mini & 0.731 & 0.756 & 3.1k \\
    \midrule
    \multirow{2}{*}{Diabetic} & \qwen{} & 0.765 & 0.749 & 3.0k \\
     & GPT-5-mini & 0.704 & 0.718 & 3.1k \\
    \midrule
    \multirow{2}{*}{Mean} & \qwen{} & 0.773 & 0.769 & 2.9k \\
     & GPT-5-mini & 0.743 & 0.746 & 3.2k \\
    \bottomrule
  \end{tabular}
  \caption{Cross-model and per-dataset breakdown at the native operating point
  (covered = non-\inconclusive{} verdicts, no score thresholding). GPT-5-mini has
  no \texttt{case\_score}, so this is the common ground for comparison; \qwen{}'s
  selective peak (\Cref{tab:headline}) is higher. Tok/prop is output tokens per
  (repository, property), a cost proxy. \qwen{} rows are from the headline FP8
  run (\Cref{tab:headline}); GPT-5-mini from a matched run.}
  \label{tab:models}
\end{table}

\paragraph{RQ4: Is the cited evidence trustworthy?}
Verdict accuracy says nothing about whether the cited evidence is sound, so we
check it directly with the independent verifier (\Cref{sec:setup}).
Across the synthetic benchmark, the verifier confirms that the cited evidence
supports the verdict for 91.3\% of the \fail{} verdicts that cite checkable
evidence (306 of 335; 95\% Wilson CI [0.88, 0.94]).
Requiring supported evidence lowers fail precision only modestly, from 0.79 to
0.73, so \vibetest{} rarely reaches a correct \fail{} for the wrong reason.
The weaker point is evidence coverage: 87 of 422 \fail{} verdicts (21\%) cite no
checkable evidence, and we count these against \vibetest{} in the 0.73 figure.
This matters because the citation is what lets a reader confirm each verdict.
The evidence verifier was run as a separate FP8 pass over the same benchmark and
model (the headline run does not store per-verdict evidence checks), so these
numbers reflect the same tester on the same cases with the verifier enabled.

\section{Auditing Synthetic Disagreements}
\label{sec:audit}

Real Kaggle appears stronger than synthetic Kaggle under the headline metrics.
One explanation is that the synthetic labels are stricter or noisier than the
evidence standard used by \vibetest{} and the verifier.
We test this directly with the LLM verifier (\Cref{sec:setup}), which
re-examines high-confidence disagreements between \qwen{} predictions and
synthetic labels and labels each as a model miss or a mismatch.
After excluding one parse error, 53 sampled disagreements were usable.
Of these, 29 were real \qwen{} misses and 24 were
property-operationalization mismatches: cases where the injected label and the
evidence-based judgment apply different but individually defensible property
boundaries.
The mismatch share is 45\% (95\% Wilson CI [33\%, 59\%]); the interval is wide
because the audit is small, but even its lower bound puts property-boundary
disagreement at roughly a third of the sampled errors.
\Cref{tab:synthetic-examples} gives representative cases of each type.

\begin{table*}[t]
  \centering
  \small
  \setlength{\tabcolsep}{4pt}
  \begin{tabularx}{\textwidth}{@{}p{2.0cm}p{1.4cm}p{2.15cm}p{2.25cm}X@{}}
    \toprule
    Type & Dataset & Repository & Property & Audit finding \\
    \midrule
    Model miss & Diabetic & \texttt{SYN-D1} & Parameters updated &
    All parameters are frozen before an initial training call; later unfreezing
    does not erase the earlier frozen training stage. \\
    Operationalization mismatch & Diabetic & \texttt{SYN-D1} & Replaceable loops &
    The synthetic label treats benign file/path loops as failures; the evidence
    rubric restricts the property to active numerical loops replaceable by
    vectorized operations. \\
    Operationalization mismatch & Diabetic & \texttt{SYN-D2} & Train-only augmentation &
    The synthetic label says \pass, but the evidence rubric flags stochastic
    transforms reused for train and validation loaders. \\
    \bottomrule
  \end{tabularx}
  \caption{Representative synthetic audit examples. Some disagreements are
  model misses, while others expose different interpretations of the property.}
  \label{tab:synthetic-examples}
\end{table*}

This audit changes how we interpret synthetic performance.
Synthetic Kaggle is valuable because it provides broad labeled coverage and
exposes real model misses.
But raw synthetic F1 is pessimistic as a model-quality estimate because it also
penalizes cases where the benchmark label and evidence standard disagree about
the property boundary.
This is especially important for natural-language ML properties, where defining
what counts as a failure is part of the task.

\section{Real-Kaggle Fail Audit}
\label{sec:real-examples}

The conservative real Kaggle audit sampled 152 predicted failures and labeled 122
as true failures and 30 as false failures, a fail precision of 0.80 (95\% Wilson
CI [0.73, 0.86]).
A full-context reaudit of those 30 conservative false failures accepted 26 as
true failures once the verifier was given relevant source context.
This supports our choice to headline conservative results while treating
full-context adjusted values as sensitivity analysis (\Cref{tab:real-sensitivity}).
Representative true failures and context-flipped cases are shown in
\Cref{tab:real-audit-examples} in the appendix: they illustrate that full source
context can change whether a predicted failure is judged supported, so
context-limited verifier results should be read conservatively.

\section{Discussion}
\label{sec:discussion}

\vibetest{} turns a natural-language ML-pipeline property into an
evidence-backed repository test.
It answers one narrow question: given a single property, does the repository
satisfy it, and what evidence shows that.
The method rests on three choices.
Because many properties are only partially observable, the three-verdict schema
keeps \inconclusive{} distinct from \fail{}.
The score allows selective prediction.
And the audit reflects a simple fact: a benchmark label encodes one
interpretation of a property, and other defensible interpretations exist.

We therefore treat the synthetic disagreement audit as a result in its own
right.
If nearly half of the high-confidence disagreements are property-boundary
mismatches, the benchmark is measuring two things at once: model behavior and
label operationalization.
This explains why synthetic raw F1 can look lower than real audited F1 without
implying that the real benchmark is simply easier.
The two settings use different standards: synthetic labels come from injected
changes, while real results come from evidence audits.
The trajectory is a familiar one.
SWE-bench's \emph{Verified} re-annotation showed that automatically generated
task labels had to be human-audited before they could be trusted
~\citep{swebench-verified}; once a property is only stated in natural language,
its labels warrant the same scrutiny.

\paragraph{Baseline interpretation.}
The baseline evidence should be read in tiers.
Reviewer modes are the strongest full-benchmark comparison because they run on
all synthetic cases; they show that less structured review prompts lag behind
\vibetest{} in selective F1.
Codex reviewer and TrainCheck are execution/format-fragile in this setting: on
the full benchmark they emit parseable committed verdicts on only 5.5\% and 0\%
of cases, so they enter \Cref{tab:headline} as native low-coverage points rather
than as selective curves.
Their fragility at scale is itself informative: producing an auditable verdict
on messy repositories is a nontrivial part of the task.
Direct prompting tests whether the gain comes from simply asking an LLM about
the property.
In the five-repository Titanic probe \vibetest{} has substantially higher
failure recall than both direct prompting and Codex review (\Cref{tab:probe}),
and on the same-checkpoint probe (\Cref{sec:setup}), where \vibetest{} and
direct prompting share one Qwen3.6-35B checkpoint, \vibetest{} still leads by
0.11 macro F1 and 0.12 failure recall.
The advantage therefore comes from the evidence-guided procedure.

\section{Related Work}

\paragraph{Repository-level agents.}
SWE-bench and SWE-bench Verified evaluate repository-scale issue resolution
~\citep{jimenez2024swebench,swebench-verified}, while Agentless shows that
carefully designed non-agentic workflows can be competitive on repair tasks
~\citep{xia2024agentless}.
RepoAudit studies autonomous repository-level code auditing
~\citep{guo2025repoaudit}.
Our setting is narrower and more evidence-driven: the agent receives a specific
natural-language property and must decide whether the repository satisfies,
violates, or fails to establish it.

\paragraph{Program analysis and ML debugging.}
Classical static analysis provides formal guarantees under explicit
specifications~\citep{cousot1979systematic,nielson1999principles}.
ML debugging systems such as UMLAUT, TrainCheck, model assertions, and TorchQL
catch important failure classes through structured checks or constraints
~\citep{schoop2021umlaut,traincheck,kang2020assertions,naik2024torchql}.
\vibetest{} is complementary because it targets properties that have not yet
been formalized.

\paragraph{Reasoning and verification with language models.}
Work on verifiers, process supervision, self-refinement, and LLM-as-judge
evaluation shows that model outputs improve when intermediate reasoning or
feedback is made explicit
~\citep{cobbe2021verifier,lightman2024process,madaan2023selfrefine,weng2023selfverification,zheng2023judge,tan2025judgebench}.
\vibetest{} applies this lesson to repository testing: the intermediate object
is an evidence-backed verdict that can itself be audited.
Our audits use a verifier in this sense, but for a different purpose than
scoring answers: a separate LLM verifier adjudicates whether a benchmark label
and an evidence-based judgment operationalize a property the same way.

\section{Limitations}
\label{sec:limitations}

The synthetic disagreement audit covers a sample of disagreements.
It estimates a major source of residual error but does not produce a corrected
synthetic benchmark.
Our confidence intervals also resample (repository, property) pairs
independently; because the properties within one repository are correlated, a
cluster bootstrap over repositories would be more conservative, and all results
come from a single \qwen{} run.
The audits are adjudicated by an LLM verifier.
Using a different model family from the tester gives cross-model independence,
but the verifier shares the tester's evidence standard and the general failure
modes of LLM judges; a human adjudication of the mismatch cases would be a
stronger check on the operationalization finding.
The real Kaggle audit is also sampled and focuses on predicted failures, so it
does not establish complete recall over all real defects.
\needsinput{a small pass-side audit (sample predicted \pass{} cases and check
for missed real violations) would bound the false-negative rate that the real
macro F1 currently cannot, and is the single result most likely to be
challenged.}
The same-checkpoint direct-property comparison is a five-repository probe on a
public OpenRouter checkpoint; scaling it to all 1{,}125 cases remains future
work.
Codex reviewer and TrainCheck run on the full synthetic benchmark but produce
almost no committed verdicts there; a cleaner containerized execution could
raise their coverage.
On real Kaggle we report no Codex or direct-property numbers, because our real
ground truth comes from audits of \vibetest{}'s own predicted failures and no
comparable audit was run for the baselines.

\section{Conclusion}

We presented \vibetest, an evidence-guided agentic tester for natural-language
ML-pipeline properties.
On synthetic Kaggle, \vibetest{} with \qwen{} outperforms reviewer-style
baselines.
On real Kaggle, conservative sampled audits estimate macro F1 around 0.93.
The main evaluation lesson is that natural-language repository testing requires
auditing both model predictions and benchmark-label operationalizations.
For ML repositories, the hard part is deciding whether repository evidence
truly establishes a violation of the property being tested.

\bibliography{vibetest}

\appendix

\section{Real Audit Examples}

\begin{table*}[t]
  \centering
  \small
  \setlength{\tabcolsep}{4pt}
  \begin{tabularx}{\textwidth}{@{}p{2.0cm}p{1.4cm}p{2.15cm}p{2.25cm}X@{}}
    \toprule
    Type & Dataset & Repository & Property & Audit finding \\
    \midrule
    True fail & Titanic & \texttt{REAL-T1} & Dataloader shuffling &
    Training and validation dataloaders are both constructed with
    \texttt{shuffle=True}. \\
    True fail & Titanic & \texttt{REAL-T2} & Metric definitions &
    ROC-AUC is computed from hard class predictions rather than probabilities
    or continuous scores. \\
    Context-flipped fail & Titanic & \texttt{REAL-T3} & Replaceable loops &
    The conservative audit lacked cited cells; full context shows a loop-based
    \texttt{predict} function despite an equivalent vectorized forward pass. \\
    Context-flipped fail & Titanic & \texttt{REAL-T4} & Device placement &
    The conservative audit lacked cited cells; full context shows mixed
    \texttt{net.to(device)}, hardcoded \texttt{.cuda()}, \texttt{.cpu()}, and
    CPU tensor usage. \\
    \bottomrule
  \end{tabularx}
  \caption{Representative real Kaggle audit examples. Full source context can
  change whether a predicted failure is judged supported, so context-limited
  verifier results should be interpreted conservatively.}
  \label{tab:real-audit-examples}
\end{table*}

\section{Artifact Map}

The current package is indexed in \texttt{supporting\_docs/}.
The main supporting files are:
\begin{itemize}
  \item \texttt{PAPER\_CLAIMS\_AND\_RESULTS.md}: frozen claims and headline numbers.
  \item \texttt{PAPER\_FIGURE\_INVENTORY.md}: figure provenance and intended use.
  \item \texttt{PAPER\_ERROR\_ANALYSIS.md}: synthetic and real audit interpretation.
  \item \texttt{PAPER\_APPENDIX\_MATERIALS.md}: prompt, threshold, and baseline checklist.
\end{itemize}

\end{document}
